\documentclass[10pt,a4paper]{amsart}
\usepackage[margin=19mm]{geometry}
\usepackage{amsmath,amssymb,mathtools}
\usepackage[scr=rsfs,cal=euler]{mathalfa}
\usepackage{xfrac}
\usepackage[unicode,hidelinks,bookmarks=false]{hyperref}
\newcommand{\ie}{i.e.\ }
\newcommand{\cf}{cf.\ }
\newcommand{\resp}{respectively\ }
\newcommand{\Subsection}{\subsection}
\providecommand{\nobreakdash}{\nobreak\hbox{-}}
\setlength{\emergencystretch}{2em}
\multlinegap0pt
\usepackage{mathtools}
\usepackage{amssymb}
\newtheorem{theorem}{Theorem}
\title{Existence of holomorphic Lie algebroid connections in higher dimensions}
\author{Indranil Biswas, Anoop Singh}
\date{Source: arXiv:2604.06003v1 (CC BY 4.0)}
\begin{document}
\maketitle

\noindent\emph{Source and licence.} Existence of holomorphic Lie algebroid connections in higher dimensions. Version arXiv:2604.06003v1 (CC BY 4.0), distributed by arXiv under the Creative Commons Attribution 4.0 International licence, http://creativecommons.org/licenses/by/4.0/, as recorded in the arXiv licence metadata for this exact version and shown on the abstract page https://arxiv.org/abs/2604.06003v1. Full licence notice: \texttt{licenses/arxiv-2604.06003-CC-BY-4.0.txt}.\par\smallskip

\noindent\textit{Editorial scope.} Counted unit: the article's theorem thm:main (source0.tex lines 356--361). The article's complete original proof of it is delivered with it (source0.tex lines 363--472, one \texttt{proof} environment of 110 lines). No mathematics is written, repaired or re-derived: the statement and the proof are the article's own lines, in the article's own notation, and no retained line is substituted or re-flowed.  Retained uncounted as the source-local context the proof needs: lines 335--338; lines 351--354.  Nothing was omitted from any retained span, so the package carries no omission record.  The retained lines cite 1 works, whose entries are transcribed verbatim from the article's own bibliography source in the e-print and delivered with short editorial labels recorded in the item's reference mapping.  No retained line contains a figure environment, an image-inclusion command or drawing code, so no image is delivered and the package declares no asset dependency.  Editorial formatting only, in addition: an AMS \texttt{amsart} preamble that restates the article's own declarations and macros used by the retained lines, namely the environments \texttt{theorem} on separate counters, together with the standard packages those lines need.  The retained environments print in this document as Theorem 1 in order of appearance, whereas the article prints the same items with the numbering of its own full document.  The complete native source of the article as distributed in the arXiv e-print for this exact version is bundled unchanged as \texttt{sources/197993/source0.tex}.
\begin{equation}
\label{eq:V_n}
(V_n, \phi_n)
\end{equation}

\begin{equation}\label{eq:s5}
0  \longrightarrow V_n  \xrightarrow{j_n}  V\big\vert_{n} \xrightarrow{\,
\psi\, }\ \mathcal{O}_X(X_n)\big\vert_{n} \longrightarrow  0.
\end{equation}

\begin{theorem}\label{thm:main}
Assume that $\mathrm{dim}_{\mathbb{C}} (X)  \geq  3$ and the anchor map $\phi  :  V  \longrightarrow 
TX$ is surjective. A holomorphic vector bundle $E$ over $X$ admits a holomorphic $(V,\,\phi)$-connection if
and only if the restriction $E\big\vert_{X_n}$ admits a holomorphic $(V_n, \phi_n)$-connection
for sufficiently large $n$, where $(V_n, \phi_n)$ is the Lie algebroid constructed in \eqref{eq:V_n}.
\end{theorem}

\begin{proof} 
Let $\mathcal{F}$ be a coherent sheaf of $\mathcal{O}_X$-modules on $X$. Then, we have the exact
sequence of coherent sheaves
\begin{equation} \label{eq:s6}
0 \longrightarrow \mathcal{F} \otimes \mathcal{O}_X(-X_n) \longrightarrow
\mathcal{F} \longrightarrow i_*\left( \mathcal{F}\big\vert_{X_n} \right) \longrightarrow 0,
\end{equation}
where $ i :  X_n  \hookrightarrow   X$ is the inclusion map. Let $\mathcal{O}_X(1)
:= i^* \mathcal{O}_{\mathbb{CP}^N} (1)$. Then, $\mathcal{O}_X(1)$ is very ample and $\mathcal{O}_X(X_n)
 \cong  \mathcal{O}_X(n)$. Since
$\dim X \geq  3$, it follows from the Serre's duality theorem \cite[Theorem 7.6, p.~243]{Ha} and Serre's theorem
\cite[Theorem 5.2, p.~228]{Ha} that for $n  \gg  0$, we have
\begin{equation} \label{eq:s7}
\mathrm{H}^i\left(X,  \mathcal{F} \otimes \mathcal{O}_X(-X_n) \right)  =  0, \quad \forall\,  i  = 1,  2.
\end{equation}
Consider the long exact sequence of cohomologies associated with the short exact sequence
\eqref{eq:s6}
$$
\mathrm{H}^1\left(X,\, \mathcal{F} \otimes \mathcal{O}_X(-X_n) \right) 
 \longrightarrow  \mathrm{H}^1(X,  \mathcal{F})  \longrightarrow
\mathrm{H}^1\left(X,  i_*\left( \mathcal{F}\big\vert_{X_n} \right)\right)  = 
\mathrm{H}^1\left(X_n, \mathcal{F}\big\vert_{X_n}\right)\,$$
$$
\longrightarrow 
\mathrm{H}^2\left(X, \mathcal{F} \otimes \mathcal{O}_X(-X_n) \right).
$$
Using \eqref{eq:s7} in this exact sequence we get an isomorphism
\begin{equation} \label{eq:s8}
\mathrm{H}^1(X, \mathcal{F})  \xrightarrow{\, \cong \,}
\mathrm{H}^1\left(X_n, \mathcal{F}\big\vert_{X_n}\right).
\end{equation}
Since $\phi$ is surjective (by assumption), we have the short exact sequence in \eqref{eq:s5}. Consider
its dual exact sequence
\begin{equation}
\label{eq:s9}
0  \longrightarrow  \mathcal{O}_X(-X_n)\big\vert_n  \longrightarrow  V^*\big\vert_n
\longrightarrow  V^*_n  \longrightarrow   0. 
\end{equation}
Tensoring the short exact sequence in \eqref{eq:s9} with $\mathrm{End} (E)\big\vert_{X_n}$ the following
short exact sequence is obtained:
\begin{equation}
\label{eq:s10}
0 \longrightarrow  \left( \mathrm{End} (E) \otimes \mathcal{O}_X(-X_n) \right)\big\vert_n
 \longrightarrow  \left( \mathrm{End} (E) \otimes V^* \right)\big\vert_n  \longrightarrow  
\mathrm{End} (E)\big\vert_n \otimes V^*_n  \longrightarrow  0.
\end{equation}

Setting $\mathcal{F}  =  \mathrm{End} (E) \otimes \mathcal{O}_X(-X_n) $ in the short exact sequence in
\eqref{eq:s6}, we get the following
\begin{equation}\label{eq:s11}
0  \longrightarrow  \mathrm{End} (E)\otimes \mathcal{O}_X(-2 X_n)   \longrightarrow  \mathrm{End} (E)\otimes
\mathcal{O}_X(-X_n)   \longrightarrow  \left( \mathrm{End} (E) \otimes \mathcal{O}_X(-X_n) \right)\big\vert_n
\longrightarrow  0.
\end{equation}
Consider the long exact sequence of cohomologies associated to \eqref{eq:s11}:
\begin{equation*}
\longrightarrow  \mathrm{H}^1\left(X,  \mathrm{End} (E)\otimes \mathcal{O}_X(- X_n) \right)   \longrightarrow 
\mathrm{H}^1\left(X_n,  \left( \mathrm{End} (E) \otimes \mathcal{O}_X(-X_n) \right)\big\vert_n \right)
 \end{equation*}
\begin{equation}\label{eq:s11b}
\longrightarrow \mathrm{H}^2\left(X,  \mathrm{End} (E)\otimes \mathcal{O}_X(-2 X_n) \right) 
\longrightarrow.
\end{equation}
Note that from \eqref{eq:s7} we have 
\begin{equation}
\label{eq:s12}
\mathrm{H}^2\left(X,  \mathrm{End} (E)\otimes \mathcal{O}_X(-2 X_n) \right)  =  0  =  
\mathrm{H}^1\left(X,  \mathrm{End} (E)\otimes \mathcal{O}_X(- X_n) \right)
\end{equation}
for all $n  \gg  0$. This and \eqref{eq:s11b} combine together to give that
\begin{equation}
\label{eq:s13}
\mathrm{H}^1\left(X_n,  \left( \mathrm{End} (E) \otimes \mathcal{O}_X(-X_n) \right)\big\vert_n \right)  =  0
\end{equation}
for all $n  \gg   0$.

Consider the long exact sequence of cohomologies for the short exact sequence in \eqref{eq:s10}:
$$
\longrightarrow  \mathrm{H}^1\left(X_n, \left( \mathrm{End} (E) \otimes
\mathcal{O}_X(-X_n) \right)\big\vert_n \right)
 \longrightarrow  \mathrm{H}^1 (X_n, \left( \mathrm{End} (E)\otimes V^* \right)\big\vert_n) 
\longrightarrow  \mathrm{H}^1 (X_n, \mathrm{End} (E)\big\vert_n \otimes V^*_n)  \longrightarrow .
$$
In view of \eqref{eq:s13}, it gives an injective homomorphism
\begin{equation}
\label{eq:s14}
\mathrm{H}^1 (X_n,  \left( \mathrm{End} (E)\otimes V^* \right)\big\vert_n)  \hookrightarrow 
\mathrm{H}^1 (X_n,  \mathrm{End} (E)\big\vert_n \otimes V^*_n).
\end{equation}
Also, from the isomorphism in \eqref{eq:s8} we have 
\begin{equation}
\label{eq:s15}
\mathrm{H}^1(X, \mathrm{End} (E)\otimes V^*)\ \xrightarrow{\, \cong \,}
\mathrm{H}^1 (X_n, \left( \mathrm{End} (E)\otimes V^* \right)\big\vert_n)
\end{equation}
for all $n  \gg  0$. Now, from \eqref{eq:s14} and \eqref{eq:s15}, the homomorphism 
$$
i^{\sharp}  :  \mathrm{H}^1(X,\, \mathrm{End} (E)\otimes V^*)  \longrightarrow
\mathrm{H}^1 (X_n, \mathrm{End} (E)\big\vert_n \otimes V^*_n)
$$
induced by the natural inclusion map $i  :  X_n  \hookrightarrow  X$ is injective, so
\begin{equation}\label{eq:s16}
i^{\sharp}  :  \mathrm{H}^1(X,  \mathrm{End} (E)\otimes V^*)  \hookrightarrow 
\mathrm{H}^1 (X_n,  \mathrm{End} (E)\big\vert_n \otimes V^*_n).
\end{equation}

Note that $i^{\sharp} (\mathrm{at}_V(E))   =  \mathrm{at}_{V_n}(E\big\vert_n)$ by the functoriality of
the Atiyah class. Thus, from \eqref{eq:s16} it follows immediately that for all $n  \gg  0$, we have $\mathrm{at}_{V_n}(E\big\vert_n)  =   0$ if and only if $\mathrm{at}_V(E)  =  0$.
This completes the proof. 
\end{proof}

\begin{thebibliography}{99}
\bibitem[Ha]{Ha} R. Hartshorne, {\it Algebraic Geometry}, Graduate Texts in Mathematics, vol. 52. Springer, New York (1977). https://doi.org/10.1007/978-1-4757-3849-0
\end{thebibliography}
\end{document}
